% Gerado por regressao/06_tabelas_latex.R (projeto Modelagem). Não editar à mão.
\begin{landscape}
\begin{table}[htbp]
\caption{Robustez R16 — ano eleitoral e pré-eleitoral separados por eleição (modelo A)}\label{tab:robustez_R16}
\centering
\footnotesize
\setlength{\tabcolsep}{3pt}
\begin{tabular}{@{}>{\raggedright\arraybackslash}p{\dimexpr\linewidth-594pt\relax}>{\centering\arraybackslash}p{60pt}>{\centering\arraybackslash}p{60pt}>{\centering\arraybackslash}p{60pt}>{\centering\arraybackslash}p{60pt}>{\centering\arraybackslash}p{60pt}>{\centering\arraybackslash}p{60pt}>{\centering\arraybackslash}p{60pt}>{\centering\arraybackslash}p{60pt}>{\centering\arraybackslash}p{60pt}@{}}
\hline
\textbf{Variável} & \textbf{\hspace{0pt}Despesa total} & \textbf{\hspace{0pt}Saúde e Saneamento} & \textbf{\hspace{0pt}Educação e Cultura} & \textbf{\hspace{0pt}Habitação e Urbanismo} & \textbf{\hspace{0pt}Assistência Social} & \textbf{\hspace{0pt}Transporte} & \textbf{\hspace{0pt}Administração} & \textbf{\hspace{0pt}Agricultura} & \textbf{\hspace{0pt}Comunicações} \\ \hline
\multicolumn{10}{@{}l}{\textit{Modelo A (ciclo eleitoral)}} \\
Ano eleitoral de 2016 & 30,935 & 24,353\textsuperscript{***} & $-$18,578 & 57,436\textsuperscript{***} & 2,994 & 17,628\textsuperscript{***} & $-$17,401\textsuperscript{**} & $-$2,341 & $-$0,389\textsuperscript{**} \\
 & (85,307) & (8,737) & (46,930) & (8,251) & (4,331) & (6,405) & (7,758) & (4,314) & (0,182) \\
Ano eleitoral de 2024 & 467,553\textsuperscript{***} & 74,464\textsuperscript{***} & 99,544\textsuperscript{***} & 180,025\textsuperscript{***} & 22,594\textsuperscript{***} & 60,162\textsuperscript{***} & 26,208\textsuperscript{***} & 2,210 & $-$1,210\textsuperscript{***} \\
 & (79,477) & (6,643) & (34,951) & (18,756) & (3,490) & (7,210) & (6,631) & (3,797) & (0,226) \\
Ano pré-eleitoral & 52,068 & $-$6,984 & 41,546\textsuperscript{*} & 15,345\textsuperscript{***} & 3,640 & 4,317 & $-$5,873 & 3,429\textsuperscript{*} & 0,427\textsuperscript{***} \\
de 2015 & (42,138) & (5,058) & (23,661) & (2,518) & (2,235) & (2,938) & (3,929) & (2,084) & (0,094) \\
Ano pré-eleitoral & $-$9,598 & $-$26,741\textsuperscript{***} & 19,599 & $-$11,842 & 2,172 & $-$10,992\textsuperscript{*} & 2,091 & $-$0,504 & 0,835\textsuperscript{***} \\
de 2019 & (68,082) & (7,210) & (32,991) & (10,661) & (3,558) & (5,662) & (5,677) & (3,224) & (0,165) \\
Ano pré-eleitoral & 392,115\textsuperscript{***} & 29,470\textsuperscript{***} & 157,564\textsuperscript{***} & 83,331\textsuperscript{***} & 20,259\textsuperscript{***} & 25,594\textsuperscript{***} & 31,265\textsuperscript{***} & 11,195\textsuperscript{*} & 0,840\textsuperscript{***} \\
de 2023 & (118,413) & (8,851) & (55,445) & (22,130) & (5,335) & (9,720) & (10,044) & (5,747) & (0,180) \\
\hline
Observações & 70.669 & 70.526 & 70.550 & 69.574 & 70.485 & 54.423 & 70.553 & 64.464 & 12.360 \\
Municípios & 5.568 & 5.568 & 5.568 & 5.564 & 5.568 & 5.205 & 5.568 & 5.436 & 1.896 \\
R\textsuperscript{2} & 0,839 & 0,754 & 0,710 & 0,304 & 0,374 & 0,192 & 0,276 & 0,177 & 0,032 \\
\hline
\end{tabular}
\fonte{Elaboração própria.}
\nota{Só modelo A: efeito fixo de município; ano eleitoral e pré-eleitoral separados por eleição (2020 absorvido pela dummy de pandemia). Erros-padrão de Driscoll-Kraay entre parênteses. \textsuperscript{***}p$<$0,01; \textsuperscript{**}p$<$0,05; \textsuperscript{*}p$<$0,1.}
\end{table}
\end{landscape}
